% A tutorial deck — how to use the HSG Beamer template.
% Compile from the repo root with:
%   pdflatex templates/demo-using-the-template.tex   (3 passes)
\documentclass[aspectratio=169]{beamer}
\usepackage{listings}
\usepackage{xcolor}
\usepackage{booktabs}

% Tell the theme where the .sty + assets live (one directory up).
\makeatletter
  \def\input@path{{../}}
\makeatother
\usetheme{HSG}

\title{Using the HSG Beamer template}
\subtitle{A 10-minute tour with examples}
\author{HSG Beamer Theme Team}
\institute{University of St.\,Gallen}
\date{}
\coursefooter{Template tutorial}
\coverimage{hsg-closing-campus.jpg}

\hsgclosingcontact
  {Template Contributors}
  {beamer-hsg@example.edu}

\hsgsetsections{Setup, Helpers, Sidebar, Tips}

% ---- listing style for showing LaTeX snippets on slides ------------------
\definecolor{HSGlst}{HTML}{F4F4F0}
\lstdefinelanguage{LaTeXmini}{
  morekeywords={titlepage,frametitle,framesubtitle,coverimage,hsgemph,
                hsgagenda,hsgsection,hsgsetsections,hsgquote,hsgquoteinverted,
                hsgstatement,hsgfullimage,hsgclosing,hsgclosingcontact,
                hsgclosingtext,hsgclosingphoto,coursefooter,documentclass,
                usetheme,title,subtitle,author,institute,date,begin,end,
                frame,item,includegraphics,column,columns,textbf,color},
  sensitive=true,
  morestring=[b]{},
  morecomment=[l]\%,
}
\lstset{
  language=LaTeXmini,
  basicstyle=\ttfamily\scriptsize,
  keywordstyle=\color{HSGgreen},
  commentstyle=\color{HSGmuted}\itshape,
  backgroundcolor=\color{HSGlst},
  frame=single,
  rulecolor=\color{HSGmuted!40},
  framerule=0.4pt,
  framesep=4pt,
  xleftmargin=4pt,
  xrightmargin=4pt,
  breaklines=true,
  showstringspaces=false,
  aboveskip=0.4em,
  belowskip=0.4em,
}

\begin{document}

% =========================================================================
% 1. TITLE
% =========================================================================
\begin{frame}[plain]
  \titlepage
\end{frame}

% =========================================================================
% 2. AGENDA
% =========================================================================
\hsgagenda{
  \item Setup: assets, compiler, three commands
  \item Helpers: every macro with one example
  \item Sidebar: declaring the deck outline
  \item Tips: traps, golden rules, where to look
}

% =========================================================================
% SECTION 1: Setup
% =========================================================================
\hsgsection{Setup}

\begin{frame}[t,fragile]
  \frametitle{Three commands to a finished deck}

  \begin{enumerate}
    \setlength\itemsep{0.5em}
    \item Copy \texttt{beamerthemeHSG.sty} and every \texttt{hsg-*}
          asset into your working directory.
    \item Write your \texttt{deck.tex} starting from the skeleton in
          \texttt{templates/skeleton.tex}.
    \item Compile three times:
  \end{enumerate}

  \vspace{0.6em}
  \begin{lstlisting}
pdflatex deck.tex
pdflatex deck.tex
pdflatex deck.tex
  \end{lstlisting}

  \vspace{0.6em}
  Three passes are needed because TikZ \texttt{remember picture}
  overlays settle on the second-and-third pass.
\end{frame}

\begin{frame}[t,fragile]
  \frametitle{Preamble: what to set once, in this order}

  \begin{lstlisting}
\documentclass[aspectratio=169]{beamer}
\usetheme{HSG}

\title{My talk}
\subtitle{Optional second line}
\author{Prof. Dr. Author Name}
\institute{Your Institute\\University of St.\,Gallen}
\date{}
\coursefooter{Course or project}
\coverimage{cover-photo.jpg}     % optional hero image

\hsgclosingcontact
  {Prof. Dr. Author Name}
  {author@unisg.ch}

\hsgsetsections{Intro, Methods, Results, Outlook}
  \end{lstlisting}
\end{frame}

% =========================================================================
% SECTION 2: Helpers
% =========================================================================
\hsgsection{Helpers}

\begin{frame}[t]
  \frametitle{The cover slide}

  Just one line in your document body:

  \fcolorbox{HSGmuted!40}{HSGsand!10}{\parbox{0.95\linewidth}{\ttfamily\scriptsize
    \textbackslash begin\{frame\}[plain]\textbackslash titlepage\textbackslash end\{frame\}}}

  \vspace{0.6em}
  The cover photo comes from \texttt{\textbackslash coverimage\{...\}}
  in the preamble. The HSG green badge frame is layered behind it.

  \vspace{0.6em}
  \fcolorbox{HSGmuted!40}{white}{%
    \parbox{0.96\linewidth}{\scriptsize
      \textbf{Tip:} the cover photo right edge clips with a slanted
      polygon. Pick a photo whose subject sits in the LEFT third so
      the clip does not crop it.}}
\end{frame}

\begin{frame}[t,fragile]
  \frametitle{Agenda --- numbered list with concrete-texture column}

  \begin{lstlisting}
\hsgagenda{
  \item Introduction
  \item Methodology
  \item Results
  \item Discussion
  \item Conclusion
}
  \end{lstlisting}

  \vspace{0.4em}
  Use 4--6 items, one or two words each. The same items should appear in
  \texttt{\textbackslash hsgsetsections} so the sidebar progress matches.
\end{frame}

\begin{frame}[t]
  \frametitle{Section advance \emph{(silent, no slide rendered)}}

  \fcolorbox{HSGmuted!40}{HSGsand!10}{\parbox{0.95\linewidth}{\ttfamily\scriptsize
    \textbackslash hsgsection\{Methods\}\quad\textcolor{HSGmuted}{\% advances the sidebar; renders no frame.}\\
    \textbackslash begin\{frame\}\\
    \quad\textbackslash frametitle\{Sub-topic\}\\
    \quad ...\\
    \textbackslash end\{frame\}}}

  \vspace{0.6em}
  Call \texttt{\textbackslash hsgsection\{Name\}} right before the first
  content slide of each section. The vertical green progress bar on the
  left fills proportionally.
\end{frame}

\begin{frame}[t]
  \frametitle{Standard content slide}

  \fcolorbox{HSGmuted!40}{HSGsand!10}{\parbox{0.95\linewidth}{\ttfamily\scriptsize
    \textbackslash begin\{frame\}\\
    \quad\textbackslash frametitle\{Results\}\\
    \quad\textbackslash framesubtitle\{Detection performance\}\\
    \quad\textbackslash begin\{itemize\}\\
    \quad\quad\textbackslash item \textbackslash hsgemph\{AUC = 0.89\} on the held-out split\\
    \quad\quad\textbackslash item False-positive rate held at 2\,\%\\
    \quad\textbackslash end\{itemize\}\\
    \textbackslash end\{frame\}}}

  \vspace{0.4em}
  \texttt{\textbackslash framesubtitle} renders as a tiny grey
  breadcrumb above the title. Keep it under 6 words.
\end{frame}

\begin{frame}[t]
  \frametitle{Two-column comparison --- native Beamer \texttt{columns}}

  \fcolorbox{HSGmuted!40}{HSGsand!10}{\parbox{0.95\linewidth}{\ttfamily\scriptsize
    \textbackslash begin\{columns\}[T]\\
    \quad\textbackslash column\{0.5\textbackslash textwidth\}\\
    \quad\quad\textbackslash textbf\{Symmetric\}\\
    \quad\quad\textbackslash begin\{itemize\}\textbackslash item Shared key\textbackslash item Fast\textbackslash end\{itemize\}\\
    \quad\textbackslash column\{0.5\textbackslash textwidth\}\\
    \quad\quad\textbackslash textbf\{Asymmetric\}\\
    \quad\quad\textbackslash begin\{itemize\}\textbackslash item Key pair\textbackslash item Slow\textbackslash end\{itemize\}\\
    \textbackslash end\{columns\}}}

  \vspace{0.4em}
  No HSG-specific helper needed; Beamer's own \texttt{columns}
  environment is the right tool.
\end{frame}

\begin{frame}[t,fragile]
  \frametitle{Statement --- one big bold sentence}

  \begin{lstlisting}
\hsgstatement{Security is a process, not a product.}
  \end{lstlisting}

  \vspace{0.4em}
  Renders at 36 pt HSG green bold on its own slide, with the standard
  footer chrome (logo + page number badge). Use sparingly --- the
  whole deck should not feel like a billboard.
\end{frame}

\hsgstatement{Security is a process, not a product.}

\begin{frame}[t,fragile]
  \frametitle{Quote slides --- two flavours}

  \textbf{Black-on-white:}
  \begin{lstlisting}
\hsgquote{The cloud is just somebody else's computer.}{Anonymous}
  \end{lstlisting}

  \vspace{0.5em}
  \textbf{White-on-green:}
  \begin{lstlisting}
\hsgquoteinverted{The best way to predict the future is to invent it.}{Alan Kay}
  \end{lstlisting}
\end{frame}

\hsgquote
  {The cloud is just somebody else's computer.}
  {Anonymous}

\begin{frame}[t,fragile]
  \frametitle{Full-bleed image}

  \begin{lstlisting}
\hsgfullimage{path/to/photo.jpg}
  \end{lstlisting}

  \vspace{0.3em}
  Image fills the whole 16$\times$9 canvas. Pre-crop to 16:9 to avoid
  axis-aligned letterboxing. A small HSG logo overlays the bottom-left.
\end{frame}

% =========================================================================
% SECTION 3: Sidebar
% =========================================================================
\hsgsection{Sidebar}

\begin{frame}[t,fragile]
  \frametitle{Declaring the deck outline}

  \begin{lstlisting}
% In the preamble:
\hsgsetsections{Intro, Methods, Results, Outlook}

% In the document body, one call per section transition:
\hsgsection{Intro}        % section 1 starts here
... slides ...
\hsgsection{Methods}      % section 2 starts here
... slides ...
  \end{lstlisting}

  \vspace{0.4em}
  The narrow vertical bar on the left of every content slide fills
  green up to the current section. Past sections appear filled, the
  current section has a small ring marker, future sections are empty.
\end{frame}

\begin{frame}[t,fragile]
  \frametitle{Closing slide}

  \begin{lstlisting}
% In the preamble:
\hsgclosingcontact
  {Prof. Dr. Author Name}
  {author@unisg.ch}

% Last line of the document body:
\hsgclosing
  \end{lstlisting}

  \vspace{0.4em}
  \texttt{\textbackslash hsgclosing} produces the
  full-bleed-campus-photo \emph{Questions?} slide, with the translucent
  contact box, the HSG logo, the EQUIS/AACSB/AMBA accreditation strip,
  and the \emph{From insight to impact} tagline. Override the headline
  with \texttt{\textbackslash hsgclosingtext\{Thank you.\}} or the photo
  with \texttt{\textbackslash hsgclosingphoto\{...\}}.
\end{frame}

% =========================================================================
% SECTION 4: Tips
% =========================================================================
\hsgsection{Tips}

\begin{frame}[t,fragile]
  \frametitle{Golden rules}

  \begin{itemize}
    \setlength\itemsep{0.55em}
    \item \hsgemph{Never use HSG green for titles or headings.} It's
          reserved for the sidebar, the page badge, and inline
          \texttt{\textbackslash hsgemph\{...\}}.
    \item \hsgemph{Three passes} of \texttt{pdflatex}. Two is not
          enough; remember-picture overlays drift on pass two.
    \item \hsgemph{Copy ALL assets} listed in the theme header. A
          missing \texttt{hsg-closing-logo-band.png} silently breaks
          the closing slide.
    \item \hsgemph{Keep slides spare.} If you find yourself shrinking
          to fit, split the slide instead.
    \item \hsgemph{Cover and closing} are hand-tuned. Do not redefine
          their templates.
  \end{itemize}
\end{frame}

\begin{frame}[t,fragile]
  \frametitle{Where to look when something breaks}

  \begin{tabular}{@{}p{4.4cm}p{8.6cm}@{}}
    \toprule
    \textbf{Symptom} & \textbf{Where to look first} \\
    \midrule
    Title slide chrome misplaced &
      Compile two more times. Remember-picture needs three passes. \\[3pt]
    Logo or photo missing &
      One of \texttt{hsg-*.png}/\texttt{hsg-*.pdf}/\texttt{hsg-*.jpg}
      not copied next to the \texttt{.tex}. \\[3pt]
    Body bullets render green &
      You added a \texttt{\textbackslash setbeamercolor} in your
      preamble. Remove it. \\[3pt]
    Sidebar progress empty &
      You didn't call \texttt{\textbackslash hsgsetsections} in the
      preamble. \\[3pt]
    Page number clipped in badge &
      Three-digit numbers fit; four-digit do not. Split your deck. \\
    \bottomrule
  \end{tabular}
\end{frame}

\begin{frame}[t,fragile]
  \frametitle{Further reading inside this repo}

  \begin{itemize}
    \setlength\itemsep{0.4em}
    \item \texttt{AGENTS.md} --- exhaustive agent-friendly reference.
    \item \texttt{templates/skeleton.tex} --- minimal compileable
          starter.
    \item \texttt{templates/skeleton-content.tex} --- one example of
          every helper.
    \item \texttt{example.tex} (this folder's parent) --- the
          ready-to-render demo deck.
    \item \texttt{OVERLEAF.md} --- how to upload the template to
          Overleaf and pick a compiler.
  \end{itemize}

  \vspace{0.6em}
  Issues, suggestions, fixes welcome at
  \hsgemph{github.com/brunobastosrodrigues/beamer-hsg}.
\end{frame}

\hsgclosing

\end{document}
